\section{Reproducibility Specification}
\label{app:hyperparameters}

This appendix specifies the configuration used to reproduce the matched
experiments.
The executable protocol binds task definitions, input and source-code hashes,
software versions, model settings, and random seeds into a contract hash.
Cached work is reused only when that complete hash matches.

\subsection{Temporal partitions}

Table~\ref{tab:temporal_contract_complete} gives inclusive forecast-origin
boundaries. The target cutoff is checked separately, so a multi-step window is
retained only when every target lies on or before that date. The initial
training partition is labeled ``adjustment'' in the execution artifacts.

\begin{table}[!htbp]
\centering
\caption{Complete chronological partitions and expected origins per site.}
\label{tab:temporal_contract_complete}
\scriptsize
\setlength{\tabcolsep}{1.8pt}
\renewcommand{\arraystretch}{1.08}
\begin{tabular}{@{}
  >{\raggedright\arraybackslash}p{0.15\linewidth}
  >{\raggedright\arraybackslash}p{0.15\linewidth}
  >{\raggedright\arraybackslash}p{0.35\linewidth}
  >{\centering\arraybackslash}p{0.19\linewidth}
  >{\centering\arraybackslash}p{0.08\linewidth}@{}}
\toprule
\textbf{Task} & \textbf{Partition} & \textbf{Origin range} &
\textbf{Declared target cutoff} & \textbf{$N$} \\
\midrule
Hourly 336--72 & Initial training & 2019-01-15--2022-12-29 & 2022-12-31 & 1,445 \\
 & Validation & 2023-01-01--2023-12-29 & 2023-12-31 & 363 \\
 & Refit & 2019-01-15--2023-12-29 & 2023-12-31 & 1,810 \\
 & Test & 2024-01-01--2024-12-28 & 2024-12-31 & 363 \\
\addlinespace
Daily 365--30 & Initial training & 2020-01-01--2022-12-02 & 2022-12-31 & 1,067 \\
 & Validation & 2023-01-01--2023-12-02 & 2023-12-31 & 336 \\
 & Refit & 2020-01-01--2023-12-02 & 2023-12-31 & 1,432 \\
 & Test & 2024-01-01--2024-12-02 & 2024-12-31 & 337 \\
\addlinespace
Monthly 12--1 & Initial training & 2020-01-31--2022-12-31 & 2022-12-31 & 36 \\
 & Validation & 2023-01-31--2023-12-31 & 2023-12-31 & 12 \\
 & Refit & 2020-01-31--2023-12-31 & 2023-12-31 & 48 \\
 & Test & 2024-01-31--2024-12-31 & 2024-12-31 & 12 \\
\addlinespace
Monthly 12--6 & Initial training & 2020-01-31--2022-07-31 & 2022-12-31 & 31 \\
 & Validation & 2023-01-31--2023-07-31 & 2023-12-31 & 7 \\
 & Refit & 2020-01-31--2023-07-31 & 2023-12-31 & 43 \\
 & Test & 2024-01-31--2024-07-31 & 2024-12-31 & 7 \\
\bottomrule
\end{tabular}
\end{table}

The cutoff column records the calendar boundary enforced by the window
constructor, not necessarily the last observed target timestamp. In particular,
the hourly test's last origin is 2024-12-28 00:00 and its 72 hourly targets end
at 2024-12-30 23:00 in the fixed local-time index; UTC timestamps for western
sites may extend into 31 December.

The 2019 observations preceding the first daily and monthly origins belong only
to the causal context. Selection scales use all available observations before
2023-01-01, and refit scales use all available observations before 2024-01-01.
Each scale is site-specific and records its observation count, minimum, maximum,
and effective range.

These partitions are rolling-origin evaluations. Parameters and scale values
are fixed before the relevant evaluation year, while the context advances and
can contain earlier observations from that year after they become available.
No parameter update occurs within validation or test.

\subsection{Model and optimization settings}

The common learned-model input is an array of shape
(batch, context length) containing only scaled GHI, with one integer site
identifier per example. The target is a vector of continuous GHI values
covering the declared forecast horizon. No learned model receives latitude, longitude, clock variables,
solar elevation, climate, cloud, precipitation, or a previously predicted
target. Table~\ref{tab:complete_hyperparameters} records the fixed model,
optimizer, seed, and ensemble settings.

\begin{table}[!htbp]
\centering
\caption{Fixed settings for all learned models.}
\label{tab:complete_hyperparameters}
\footnotesize
\setlength{\tabcolsep}{3pt}
\renewcommand{\arraystretch}{1.10}
\begin{tabular}{@{}
  >{\raggedright\arraybackslash}p{0.27\linewidth}
  >{\raggedright\arraybackslash}p{0.66\linewidth}@{}}
\toprule
\textbf{Item} & \textbf{Value} \\
\midrule
Seeds & 11, 23, 42, 67, 89 for daily/monthly; 42 for the official hourly
protocol and its DilatedRNN extension \\
Neural objective & Mean squared error on site-wise min--max normalized GHI \\
Optimizer & Adam; learning rate $10^{-3}$; weight decay $10^{-5}$ \\
Batch/selection & Hourly and daily: batch 128, 30 epochs maximum, patience 5;
monthly: batch 32, 300 epochs maximum, patience 30 \\
Determinism & Python, NumPy, and PyTorch generators fixed; deterministic
PyTorch algorithms requested; one PyTorch CPU thread \\
Site representation & Learned embedding of width 8 in TimesNet, DilatedRNN,
and LSTM; one-hot vector for XGBoost \\
TimesNet & $d_{\mathrm{model}}=8$, $d_{\mathrm{ff}}=16$, one TimesBlock,
top-$k=3$, two Inception kernels, dropout 0.1, direct linear temporal head;
sample-specific FFT period selection and per-window standardization with
$\epsilon=10^{-5}$ \\
DilatedRNN & Vanilla tanh-RNN cells with recurrent dilations $(1,2,4)$,
16 units in each layer, dense width 16, dropout 0, direct output head \\
LSTM & One recurrent layer, hidden width 32, direct output head \\
XGBoost & Squared-error objective, 150 trees, depth 3, learning rate 0.05,
subsample 0.8, column subsample 0.8, histogram tree method, one thread;
multi-output-tree strategy when the output length exceeds one \\
Refit rule & Fresh initialization under the same seed and exactly the epoch
count selected using 2023 validation data; XGBoost is independently fitted with its fixed budget \\
Ensemble & Arithmetic mean of five seed forecasts for daily/monthly learned
models; no hourly ensemble \\
\bottomrule
\end{tabular}
\end{table}

Hourly daily-profile persistence repeats the last 24 observed hourly values
cyclically over the forecast horizon. The annual seasonal na\"ive reference
uses the value at the same local hour in the preceding year. For daily and
monthly tasks, persistence repeats the last observed aggregate throughout
the forecast horizon, whereas annual seasonal na\"ive uses the same civil
day or month from the preceding year, with 29 February mapped to 28 February
when needed. Daily and monthly climatology uses a fixed mean for the target
calendar day or month, computed from observations before the relevant
selection or test cutoff. A global mean computed before the cutoff is used
when no observations match the calendar day or month. All baseline and learned
forecasts are clipped to zero if negative. Hourly outputs are also set to zero
when solar elevation is nonpositive, evaluated at minute 30 of each interval.

\subsection{Evaluation outputs and reproducibility records}

The retained corrected archive is identified as
\path{avaliacao_multirresolucao_corrigida_v2} under the project's
\path{resultados/} directory. Its task directories are
\path{hourly_72_extension}, \path{daily_30}, \path{monthly_1}, and
\path{monthly_6}. This identifier specifies the source of the reported
comparisons; it does not imply that the full forecast bundle is currently
available in the public repository.

For every declared horizon, the pipeline writes cumulative-prefix and
exact-lead metrics. It stores raw and post-processed forecasts, metrics by site,
equal-weight macro metrics, site-paired TimesNet--DilatedRNN differences,
selected epoch counts, training histories, and the two sets of scales fitted before the respective cutoffs.
Daily and monthly task directories additionally store variability across the
five seeds. The principal machine-readable files across task directories are:
\begin{itemize}
  \item \path{contrato_execucao.json} and
        \path{hiperparametros_e_metodo.json};
  \item \path{protocolo_temporal.csv},
        \path{escalas_minmax_pre_corte.csv}, and
        \path{epocas_selecionadas.csv};
  \item \path{previsoes_validacao.csv.gz} and
        \path{previsoes_teste.csv.gz};
  \item \path{metricas_por_localidade.csv},
        \path{metricas_macro.csv}, and
        \path{variabilidade_sementes.csv} (daily/monthly only);
  \item \path{comparacao_pareada_localidades.csv},
        \path{comparacao_pareada_resumo.csv}, and
        \path{manifesto_artefatos.json}.
\end{itemize}
The contract records Python, NumPy, pandas, PyTorch, scikit-learn, XGBoost, and
joblib versions together with SHA-256 hashes of source code and input files.
Smoke mode reduces sites, origins, seeds, epochs, and trees and is marked
nonpublishable in every status artifact.

\subsection{Geographic and meteorological provenance}

The geographic CSV retains all factory-coordinate provenance, raster-cell
row and column, cell-center coordinates, sampling distance, class code and
descriptions, confidence, source citation, file paths, and the MD5 values in
Table~\ref{tab:context_hashes}. MD5 is used only as a reproducibility checksum,
not as a security mechanism.

\begin{table}[!htbp]
\centering
\caption{Integrity records for the external contextual layer.}
\label{tab:context_hashes}
\scriptsize
\setlength{\tabcolsep}{2.5pt}
\renewcommand{\arraystretch}{1.10}
\begin{tabular}{@{}
  >{\raggedright\arraybackslash}p{0.31\linewidth}
  >{\raggedright\arraybackslash}p{0.62\linewidth}@{}}
\toprule
\textbf{Asset} & \textbf{Role and retained checksum} \\
\midrule
Beck present 1-km class raster & Site assignment; MD5
\texttt{fe8f23532bfb5923\allowbreak{}1a072398f0d9a7cb} \\
Beck present 1-km confidence raster & Assignment confidence; MD5
\texttt{a87f331e406e4ef0\allowbreak{}2114e1350eb2dd01} \\
Beck present 0.5-degree raster & Americas-map background only; MD5
\texttt{78d09475fa4b0dbb\allowbreak{}004b6178cba47b4a} \\
Beck class legend & Codes, descriptions, and RGB colors; MD5
\texttt{f0632ca1b3b475d3\allowbreak{}de5a143c4ee1928f} \\
\path{contexto_geografico.csv} & Auditable sampled table; SHA-256
\texttt{80b6e003a326d614\allowbreak{}2df03766c15a5dbe\allowbreak{}%
28c7c1dc0be3c195\allowbreak{}df07a60b764f7f41} \\
\path{contexto_meteorologico_2024.csv} & Post-hoc NASA POWER context;
SHA-256
\texttt{7fbf97076193f246\allowbreak{}7e0092cf6b0d2c9a\allowbreak{}%
00ea19c84f354b28\allowbreak{}f1736f5bba95f4b5} \\
\bottomrule
\end{tabular}
\end{table}

The climate source is the 1980--2016 nominal 1-km K\"oppen--Geiger product of
Beck et al. \cite{beck2018}. Class confidence records agreement among the
12 ensemble maps used to derive that classification. The meteorological artifact contains 3,660
location--day rows in local solar time and is accompanied by a JSON manifest
that records the NASA POWER query URL and SHA-256 of every cached response
\cite{nasapowerdaily}. These data are used only to contextualize forecast
errors after evaluation. A day is classified as adverse if precipitation is
at least 1~mm/day or cloud amount is at least 80\%. The manifest explicitly
records that these variables are not model inputs.
